\clearpage
\VerdierPhysicalPageAnchor{0138}%
\begingroup
\small
\setlength{\emergencystretch}{2em}%
\setlength{\parskip}{0pt}%
\setlength{\abovedisplayskip}{1pt plus 1pt minus 1pt}%
\setlength{\belowdisplayskip}{1pt plus 1pt minus 1pt}%

\begingroup\itshape
\VerdierSemanticAnchor{ch2-2-2-4-b}
\noindent b) Lorsque la catégorie $\mathcal C$ est additive, la catégorie
$\mathcal C(S^{-1})$ est additive et le foncteur $Q$ est additif.

\VerdierSemanticAnchor{ch2-2-2-4-c}
\par\noindent c) Soit $\widehat S$ l'ensemble des morphismes de $\mathcal C$
qui sont transformés par $Q$ en isomorphismes de $\mathcal C(S^{-1})$.
\VerdierSemanticAnchor{ch2-2-2-4-c-saturation}%
L'ensemble $\widehat S$ est un système multiplicatif saturé de $\mathcal C$.
Tout système multiplicatif saturé de $\mathcal C$ qui contient $S$ contient
$\widehat S$.
\endgroup

\VerdierSemanticAnchor{ch2-2-2-4-proof-b}
\par\noindent La première assertion résulte immédiatement de la propriété
universelle de $\mathcal C(S^{-1})$ : la catégorie
$\bigl(\mathcal C^\circ(S^{-1})\bigr)^\circ$ est une solution de ce problème.
La deuxième assertion résulte immédiatement des assertions b) et c) du
théorème (2.2.2).
\VerdierSemanticAnchor{ch2-2-2-4-proof-c-characterization}%
Il résulte immédiatement de la description des morphismes de
$\mathcal C(S^{-1})$ qu'un morphisme $f$ de $\mathcal C$ est transformé par
$Q$ en isomorphisme si et seulement s'il existe deux morphismes de
$\mathcal C$, $g$ et $h$, tels que $gf$ et $fh$ soient définis et
appartiennent à $S$.
\VerdierSemanticAnchor{ch2-2-2-4-proof-c-sm2-sm3}%
On en déduit immédiatement que l'ensemble $\widehat S$ possède les
propriétés SM2) et SM3).
\VerdierSemanticAnchor{ch2-2-2-4-proof-c-sm1-sm4}%
De plus, il résulte de la définition de $\widehat S$ que cet ensemble possède
les propriétés SM1) et SM4).
\VerdierSemanticAnchor{ch2-2-2-4-proof-c-conclusion}%
L'ensemble $\widehat S$ est donc un système multiplicatif saturé de
$\mathcal C$. La dernière assertion de (2.2.4) est triviale.

\VerdierSemanticAnchor{ch2-2-2-5}
\smallskip
\noindent\textbf{2.2.5.} Soient $\mathcal C$ une catégorie et $S$ un système
multiplicatif de $\mathcal C$.
\VerdierSemanticAnchor{ch2-2-2-5-localized-category}%
La catégorie $\mathcal C(S^{-1})$ est appelée la \emph{catégorie localisée}
de $\mathcal C$ par rapport au système $S$.
\VerdierSemanticAnchor{ch2-2-2-5-localization-functor}%
Le foncteur $Q\colon\mathcal C\longrightarrow\mathcal C(S^{-1})$ est appelé
le \emph{foncteur de localisation}.
\VerdierSemanticAnchor{ch2-2-2-5-saturation}%
Le système multiplicatif $\widehat S$ est appelé le \emph{saturé} du système
$S$.

\VerdierSemanticAnchor{ch2-2-2-6}
\smallskip
\noindent\textbf{Théorème 2.2.6.}\quad
\begingroup\itshape
\VerdierSemanticAnchor{ch2-2-2-6-setup}%
Soient $\mathcal D$ une catégorie triangulée et $S$ un système multiplicatif
de $\mathcal D$ (2.1.4). Désignons par $\mathcal D(S^{-1})$ la catégorie
localisée de $\mathcal D$ et
$Q\colon\mathcal D\longrightarrow\mathcal D(S^{-1})$ le foncteur de
localisation.

\VerdierSemanticAnchor{ch2-2-2-6-a}
\par\noindent a) Il existe sur $\mathcal D(S^{-1})$ une et une seule structure
de $\mathbb Z$-catégorie stricte telle que $Q$ soit un $\mathbb Z$-foncteur
strict.

\VerdierSemanticAnchor{ch2-2-2-6-b}
\par\noindent b) Il existe sur $\mathcal D(S^{-1})$ une et une seule structure
triangulée telle que $Q$ soit un foncteur exact. Les triangles distingués
pour cette structure sont les triangles isomorphes aux images par $Q$ des
triangles distingués de $\mathcal D$.

\VerdierSemanticAnchor{ch2-2-2-6-c-factorization}
\par\noindent c) Tout foncteur exact (resp. cohomologique)
$F\colon\mathcal D\longrightarrow\mathcal D'$ dans une catégorie triangulée
(resp. abélienne) qui transforme les morphismes de $S$ en isomorphismes se
factorise d'une manière unique par
$Q\colon\mathcal D\longrightarrow\mathcal D(S^{-1})$ en un foncteur exact
(resp. cohomologique)
$G\colon\mathcal D(S^{-1})\longrightarrow\mathcal D'$.
\VerdierSemanticAnchor{ch2-2-2-6-c-morphism-lifting}%
Soient $F_1,F_2\colon\mathcal D\rightrightarrows\mathcal D'$ deux foncteurs
exacts (resp. cohomologiques) qui transforment les morphismes de $S$ en
isomorphismes ;
$G_1,G_2\colon\mathcal D(S^{-1})\rightrightarrows\mathcal D'$ leurs
factorisations respectives et
$m\colon F_1\longrightarrow F_2$ un morphisme de foncteurs exacts
(resp. cohomologiques). Il existe un et un seul morphisme de foncteurs
$p\colon G_1\longrightarrow G_2$ tel que $m=p\star Q$.
\endgroup

\VerdierSemanticAnchor{ch2-2-2-6-proof-a}
\par\noindent La première assertion est une conséquence immédiate du théorème
(2.2.2) et de SM5).
\VerdierSemanticAnchor{ch2-2-2-6-proof-c-p139-boundary}%
La troisième assertion est une conséquence immédiate de la

\VerdierSemanticAnchor{ch2-2-2-6-p139-boundary}
\endgroup
